\section{Evaluation Design}\label{sec:evaluation-design}

\begin{strip}
    \centering
    \captionof{table}{Evaluation stages, inputs, and outputs.}
    \label{tab:experimental-setup}
    \scriptsize
    \setlength{\tabcolsep}{3pt}
    \renewcommand{\arraystretch}{1.15}
    \begin{tabular}{p{0.12\textwidth} p{0.20\textwidth} p{0.40\textwidth} p{0.20\textwidth}}
        \hline
        \textbf{Stage} & \textbf{Input}                                                                     & \textbf{Processing or model}                                                                                                          & \textbf{Output}                                                                                                 \\
        \hline
        Pipeline       & Four authorisation exports.                                                        & Source-specific parsing, field normalisation, reference matching, and separation into final, review, and discarded rows.              & 1,530 final rows, 10,160 review rows, and 183,507 discarded rows in the joined modelling data.                  \\
        Classifier     & 195,197 joined rows: 11,690 positive final/review rows and 183,507 discarded rows. & TF-IDF text features with balanced logistic regression. An 80/20 grouped holdout is split by dataset, microbe, pathogen, and product. & 155,041 training rows and 40,156 holdout rows; accuracy, macro F1, ROC AUC, and average precision of $1.000$.   \\
        Ranking        & 1,530 accepted Portuguese final rows and 595 unique context--BCA interactions.     & One BCA edge is held out from each of 109 complete contexts containing at least two BCAs.                                             & Hybrid direct-support and collaborative ranking, with inferred-only results in Table~\ref{tab:ranking-results}. \\
        \hline
    \end{tabular}
\end{strip}

We evaluate the pipeline, classifier, and ranker separately.
Pipeline checks confirm that the parsers produce the expected schema, required columns exist, target expansion produces auditable rows, and each data set produces the expected \textit{final}, \textit{review}, and \textit{discarded} files.

The classifier uses the generated classes as weak labels.
Rows in \textit{final} and \textit{review} receive the positive label because they contain BCA evidence or need BCA related review.
The rows in \textit{discarded} receive the
negative label because the pipeline did not find a BCA signal. Only rows with complete scientific crop and pathogen
pairs enter training. For each row, the model concatenates labelled values from the microbial active substance,
scientific crop, scientific pathogen, product, product function, dataset, and pathogen taxon specificity fields.
Common crop and pathogen names are excluded because their language varies across datasets.

The TF-IDF vectorizer lowercases this text and represents both unigrams and bigrams, meaning individual terms and
adjacent two-term phrases. Its weight is higher for a term that is frequent in one row but uncommon across the
corpus, which lets the classifier distinguish informative names and phrases from broadly repeated terms. Terms seen
in only one row are removed with $\mathit{min\_df}=2$, and the vocabulary is capped at 60,000 features. Balanced
logistic regression uses the liblinear solver with 1,000 iterations. Reference-match flags, extracted genera, review
and discard reasons, and purpose text are excluded to reduce direct leakage from the pipeline decisions.

We use one \texttt{GroupShuffleSplit} with a 20\% holdout, random state 42. The group key is the lowercased
combination of dataset, microbial value, common and scientific pathogen names, and product name. Thus repeated or
closely related authorisation rows with the same key stay on one side of the split instead of appearing in both
training and holdout. The stored split contains 155,041 training rows and 40,156 holdout rows. These metrics measure
reproduction of labels derived from authorisation data.

Table~\ref{tab:classifier-metrics} reports the classifier metrics on the holdout set.

\begin{table}[H]
    \caption{Classifier metrics on the grouped holdout set ($n=40{,}156$).}
    \label{tab:classifier-metrics}
    \centering
    \small
    \begin{tabular}{lr}
        \hline
        \textbf{Metric}   & \textbf{Value} \\
        \hline
        Accuracy          & 1.000          \\
        Macro F1          & 1.000          \\
        ROC AUC           & 1.000          \\
        Average precision & 1.000          \\
        \hline
    \end{tabular}
\end{table}

Accuracy is the proportion of holdout rows assigned the correct class at a probability threshold of $0.5$.
Macro F1 is the mean F1 score for the positive and negative classes, giving both classes equal importance despite their
unequal frequencies. ROC AUC measures how well the predicted probabilities separate the two classes across thresholds.
Average precision summarises the precision--recall trade-off for the positive class.

All four metrics are $1.000$ because the two generated classes are strongly separated by the information in the
authorisation rows. Final and review rows contain detected BCA evidence, while discarded rows do not contain an
accepted BCA signal. The classifier receives the microbial active substance, product name, and product-function fields,
which preserve this distinction in the input text. The grouped split prevents repeated dataset--microbe--pathogen--product
groups from appearing on both sides, but it does not make the labels independent biological observations. The values
therefore show reproduction of the pipeline's weak labels, not biological efficacy, safety, or legal suitability.

For ranking, a context is one unique crop--pathogen query, identified by its scientific crop and scientific pathogen
names. It is not a source row or a product. A context may contain several standardized BCA values and repeated
authorisation rows. We aggregate repeated rows into one context--BCA interaction while retaining product and support
counts as evidence. We evaluate the 109 complete contexts with at least two distinct observed BCAs, holding out one
deterministic context--BCA edge from each. We remove the edge before
building features, then rank the held-out BCA among the remaining candidate catalogue. For reproducibility, contexts
are ordered by their scientific crop and pathogen names and the lexicographically last standardized BCA value is
selected as the held-out edge. Unavailable targets are counted as misses; 107 of 109 inferred-only cases remain
covered (98.2\%). We report results for all candidates and for an inferred-only view, which removes candidates that
still have a known context match and corresponds to the query interface. The all-candidate hybrid result is reported
below for completeness.

We measure ranking quality with HitRate@1, HitRate@5, HitRate@10, mean reciprocal rank (MRR), and normalised
discounted cumulative gain (NDCG), following standard cumulative-gain evaluation for ranked retrieval.
We report candidate coverage separately because the ranker cannot recover a BCA
that never appears elsewhere in the interaction catalogue in this warm start experiment. The baselines rank BCAs
by pathogen support and total support.

This protocol tests retrieval of observed authorisation associations. It does not establish biological efficacy,
safety, suitability for an untreated crop and pathogen pair, or legal authorisation for a new use.
